\scanpage{80}
on veut faire abstraction des questions de
représentabilité (sauf pour les $\Gamma^{\nu_i}(X_i)$)
\[
\Gamma^{\nu_1}(X_1)\otimes\cdots\otimes\Gamma^{\nu_r}(X_r)\longleftarrow\Gamma^\nu(X_1,\ldots,X_r)
\]
on aura un hom fonctoriel en
$M$ ds $P_1$ :
\[
P_{\nu}^{\unsure{P|\circ|\mathfrak{g}}}(X_1,\ldots,X_r;M)\longleftarrow P_{1,\ldots,1}(\Gamma^{\nu_1}X_1,\ldots,\Gamma^{\nu_r}X_r;M)
\]
par composition avec les morphismes canoniques
dans $P$
\[
X_1\xrightarrow{\gamma_{\nu_1}}\Gamma^{\nu_1}X_1,\ \ldots,\ X_r\xrightarrow{\gamma_{\nu_r}}\Gamma^{\nu_r}(X_r)
\]
et on suppose que c'est \inserted{tjrs} un iso.
Moyennant ceci, on a
\[
\Gamma^n(X_{\deleted{\unsure{i}}1}\times\cdots\times X_r)\xleftarrow{\sim}
\sum_{\substack{\nu=(\nu_i)\in\mathbb N^r\\\sum\nu_i=n}}
\underbrace{\Gamma^{\nu_1}(X_1)\otimes\cdots\otimes\Gamma^{\nu_r}(X_r)}_{\unsure{\Vert|=}\ \Gamma^{\nu_1,\ldots,\nu_r}(X_1,\ldots,X_r)}
\]

[ \uline{NB} Il suffirait de l'établir \unsure{avec les|sous les}
hypothèses relatives $r=2$ -- pour $r=2$ d'abord,
et le déduire ensuite pour $r$ \unsure{qcque|quelc.}.

En particulier, on trouve, \deleted{que} si on
considère le produit $X\times0$, pour $\mathbb{1}=\Gamma^0(0)$
\deleted{$\Gamma^1(X\times\unsure{1|0})\simeq$}
\[
X\otimes\mathbb{1}\simeq\Gamma^1(X)\otimes\Gamma^0(0)\simeq\Gamma^{1,0}(X,0)\simeq X
\]
\inserted{Car $P_{1,0}^{\unsure{*|P|\mathfrak{g}}}(X;0;M)\simeq P_1(X,M)$}
et de $=$ $\mathbb{1}\otimes X\simeq X$ \deleted{\unsure{$\mathbb{1}\otimes\mathbb{1}$ aussi}}
\deleted{\unsure{$X\otimes\mathbb{1}$}}

% slanted
\margin{\unsure{trouver} un \unsure{critère} simple \unsure{en genre}}
On suppose qu'il y a des produits tensoriels
stricts d'un \unsure{nb|nombre} fini de facteurs, et que
l'on a des $\Gamma^n X$ satisfaisant à la condition
plus haut.
% archive sheet number: 79
